\documentclass[11pt,a4paper]{article}
\usepackage[T1]{fontenc}
\usepackage{lmodern}
\usepackage[margin=27mm]{geometry}
\usepackage{amsmath,amssymb,amsthm,booktabs,microtype}
\usepackage[hidelinks]{hyperref}
\newtheorem{theorem}{Theorem}
\newtheorem{lemma}[theorem]{Lemma}
\newtheorem{corollary}[theorem]{Corollary}
\newcommand{\Ostar}{O^{*}}
\emergencystretch=2em
\clubpenalty=10000
\widowpenalty=10000
\title{Near-Twin Neighborhoods as a Structural Exit\\
\large A deterministic bound from three-state pair compression}
\author{}
\date{8 October 2026}
\begin{document}
\maketitle

\begin{abstract}
Pairs with almost identical signed neighborhoods are a concrete source of
correlated gap functions. We prove a structural algorithm for this case:
if $m$ disjoint pairs disagree at an average of at most $r$ other paired
variables, Max
2-SAT can be solved exactly in polynomial space and time
$\Ostar(2^n(3/4)^{m/(r+1)})$. Thus a linear number of such pairs and
constant $r$ give a fixed exponential base below two, independently of
clause density. The proof chooses compatible pairs, assigns the remaining
variables, and removes one of the four states of each retained pair by
domination. This extends the exact interaction-class rule in the earlier
notes to local errors. It does not prove that small total gap necessarily
forces near-twin pairs, and makes no literature novelty claim.
\end{abstract}

\section{The neighborhood property}

Write the objective, after multiplying by four and omitting a constant,
as
\begin{equation}
 f(a)=\sum_i h_i a_i-\sum_{i<j}J_{ij}a_i a_j,
 \qquad a_i\in\{-1,+1\}.                           \label{eq:score}
\end{equation}
For a single binary clause on $ij$, $J_{ij}$ is the product of its two
literal signs. Initial unit clauses add twice their signs to $h$.
The theorem holds for general rational quadratic coefficients of
polynomial encoding length as well; the simple-clause setting has
$J_{ij}\in\{-1,0,+1\}$.

A pair $P=\{u,v\}$, with an orientation $\varepsilon\in\{-1,+1\}$,
has discrepancy set
\begin{equation}
 D(P,\varepsilon)=
 \{w\notin\{u,v\}:J_{uw}\ne\varepsilon J_{vw}\}.      \label{eq:D}
\end{equation}
Call it an $r$-near-twin pair when $|D(P,\varepsilon)|\le r$.
The case $\varepsilon=-1$ allows negating one variable before comparing
the neighborhoods. Fields $h_u,h_v$ may be completely different.
Neither the degree nor the size of the common neighborhood is bounded.

For simple clauses, discrepancy can be read from co-occurrence counts.
Let $s_{uv}$ be the number of variables adjacent to exactly one of $u,v$,
excluding the two centers. Let $o_{uv}^{\varepsilon}$ count common
neighbors where the interaction signs disagree after orientation.
Then
\begin{align}
 |D(P,\varepsilon)|&=s_{uv}+o_{uv}^{\varepsilon},     \label{eq:dist}\\
 \sum_{w\notin P}(J_{uw}-\varepsilon J_{vw})^2
 &=s_{uv}+4o_{uv}^{\varepsilon}
  =d_u^{\rm ext}+d_v^{\rm ext}-2\varepsilon C_{uv}, \label{eq:gram}\\
 C_{uv}&=\sum_{w\notin P}J_{uw}J_{vw}.
\end{align}
Here $d_u^{\rm ext}$ excludes a possible edge to $v$. Thus a near-twin
test is an explicit signed-neighborhood counting test.

\section{The local three-state rule}

\begin{lemma}[Pair domination]\label{lem:pair}
Suppose $J_{uw}=\varepsilon J_{vw}$ for every other remaining variable
$w$. For any fixed assignment of variables already removed, an optimum
of the residual soft objective uses one of three explicitly determined
assignments of $u,v$.
\end{lemma}
\begin{proof}
Let $H_u,H_v$ be the residual fields, including all effects of the fixed
assignment, and set $z_u=a_u$, $z_v=\varepsilon a_v$.
The pair interacts with every other variable only through $z_u+z_v$.
Its internal contribution is $-\varepsilon J_{uv}z_uz_v$.
The two mixed states $(z_u,z_v)=(+1,-1)$ and $(-1,+1)$ have the same
external and internal contributions. Their field contributions are
$H_u-\varepsilon H_v$ and $-H_u+\varepsilon H_v$.
Keep the first when $H_u\ge\varepsilon H_v$, and the second otherwise.
Together with $(+1,+1)$ and $(-1,-1)$, these are the three states.
The discarded state is no better for any assignment of the other
remaining variables. Ties may be broken consistently.
\end{proof}

This is an actual domination rule, not an assumption about independence
or about how many local optima exist. The choice of mixed state must use
the \emph{residual} fields: it can change when an exceptional neighbor is
assigned.

\section{Making many pairs compatible}

Suppose $P_1,\ldots,P_m$ are disjoint near-twin pairs, with specified
orientations. Form a graph $H$ on these pairs. Join $i,j$ if
\[
 D(P_i,\varepsilon_i)\cap P_j\ne\varnothing
 \quad\text{or}\quad
 D(P_j,\varepsilon_j)\cap P_i\ne\varnothing.
\]
An edge says that retaining both pairs would prevent at least one from
being an exact twin pair. Discrepancies at variables outside all pairs
do not create edges: those variables will be assigned before compression.

\begin{lemma}[Compatible family]\label{lem:compatible}
Let $U=\bigcup_iP_i$ and
$R=\sum_i|D(P_i,\varepsilon_i)\cap U|$ be the total number of
discrepancies at other paired variables. Then $H$ has a polynomial-time
computable independent set of size at least
$m^2/(m+R)$. In particular, $R\le rm$ gives size at least $m/(r+1)$.
\end{lemma}
\begin{proof}
Orient the variables inside pairs $i,j$ according to their signs
$\varepsilon_i,\varepsilon_j$. Their interactions form a $2\times2$
matrix. Pair $i$'s discrepancies count unequal entries in its columns;
pair $j$'s discrepancies count unequal entries in its rows. If these
four entries are not all equal, at least two of the four comparisons
are unequal: three equalities around a four-cycle imply the fourth.
Consequently, every edge of $H$ accounts for at least two discrepancies,
and $2|E(H)|\le R$.

In a uniformly random ordering of $V(H)$, retain every vertex that
precedes all its neighbors. These vertices form an independent set,
of expected size
\[
 \sum_i\frac1{\deg_H(i)+1}
 \ge\frac{m^2}{m+2|E(H)|}\ge\frac{m^2}{m+R}.
\]
The first inequality is Cauchy--Schwarz. This construction is
deterministic by conditional expectations: at each position, choose a
vertex whose placement does not decrease the conditional expected
size. The expectation is computable in polynomial time. A not-yet-placed
vertex contributes zero if it already has a placed neighbor, and
otherwise contributes $1/(\deg_H(i)+1)$; placed retained vertices
contribute one.

If each individual discrepancy count is at most an integer $r$, a
simpler certificate is available: every induced subgraph has average
degree at most $r$ by the same counting argument, so reverse greedy
coloring uses at most $r+1$ colors. Its largest color class suffices.
\end{proof}

More generally, any independent set of $s$ pairs in $H$ can be used.
For example, a bipartite $H$ permits $s\ge m/2$ even if individual
discrepancy counts grow without bound. A star permits keeping all its
leaves. Thus bounded discrepancies are one sufficient certificate,
not the only useful form of this structural condition.

\begin{theorem}[Near-twin algorithm]\label{thm:main}
Given $m$ disjoint oriented pairs in a formula on $n$ variables,
with total internal discrepancy count $R\le rm$,
there is a deterministic exact algorithm using polynomial space and
time
\begin{equation}
 \Ostar\left(2^n(3/4)^{m/(r+1)}\right).             \label{eq:time}
\end{equation}
For fixed $\beta>0$ and $r\ge0$, at least $\beta n$ disjoint pairs
with this average discrepancy bound give a fixed base below two:
\[
 c_{\beta,r}=2(3/4)^{\beta/(r+1)}<2.
\]
There is no restriction on clause density or initial literal gaps.
\end{theorem}
\begin{proof}
Use Lemma~\ref{lem:compatible} to choose $s\ge m/(r+1)$ compatible
pairs. Enumerate assignments of the $n-2s$ variables outside them.
All discrepancies of each retained pair have now been assigned, so
Lemma~\ref{lem:pair} applies to each pair. Their three-state reductions
are simultaneously sound: replacing a discarded mixed state preserves
all interaction contributions and never decreases the objective.

For each outside assignment, enumerate the $3^s$ retained pair states
and evaluate the objective. The total number of candidates is
\[
 2^{n-2s}3^s=2^n(3/4)^s.
\]
The outside assignments and pair states can both be traversed depth
first with polynomial storage and polynomial work per candidate.
\end{proof}

The theorem includes the simpler full-neighborhood class with $m$
disjoint $r$-near-twin pairs. For that class the pairs need not be supplied.
Construct a compatibility graph on the original variables: edge $uv$
exists when at least one orientation satisfies~\eqref{eq:D} with at most
$r$ discrepancies. Record a qualifying orientation. A maximum matching
finds the largest disjoint collection in polynomial time. Alternatively,
a simple greedy maximal matching loses at most a factor two in $m$,
still giving a fixed exponential saving without a matching subroutine.
Once a collection is given, ignoring discrepancies outside its union
can improve the bound further; finding the best collection under that
weaker restricted test is not claimed to be easy.

If all variables are covered by pairs ($m=n/2$), illustrative bases are
\[
 r=0:\ 1.732051,\qquad
 r=1:\ 1.861210,\qquad
 r=2:\ 1.906369.
\]
These are conservative worst-case bounds. The actual discrepancy graph
often permits retaining substantially more pairs.

\section{Why this relates to small gaps}

Let $A$ be an assigned set not containing $u,v$. Their effective gaps
are $\Delta_i(A)=h_i-\sum_{w\in A}J_{iw}a_w$. Therefore
\begin{equation}
 \Delta_u(A)-\varepsilon\Delta_v(A)
 =h_u-\varepsilon h_v
  -\sum_{w\in A}(J_{uw}-\varepsilon J_{vw})a_w.       \label{eq:coupled}
\end{equation}
For simple clauses and an $r$-near-twin pair,
\[
 |\Delta_u(A)-\varepsilon\Delta_v(A)|
 \le |h_u-\varepsilon h_v|+2r.
\]
Thus when the initial fields also agree after orientation, their gaps
remain close under \emph{every} common partial assignment, however large
their degrees are. If both are exposed to a uniform assignment on $A$,
the variance of the difference of their centered gaps is exactly
$\sum_{w\in A}(J_{uw}-\varepsilon J_{vw})^2\le4r$.

The extreme example is a family with the same exposed neighborhood and
the same field. Their gaps are copies of one signed sum. When its degree
$d$ is even and its field is zero, all those gaps vanish together on
$\binom d{d/2}$ of the $2^d$ assignments of the exposed neighborhood.
The fraction is $\Theta(d^{-1/2})$, regardless of how many variables
share the gap function. This illustrates how simultaneous small gaps
can fail to behave as many independent events. It is a statement about
that family's gap functions, not by itself a bound for every earlier
gap in a full branching order.

The structural algorithm does not need the fields to agree. If they
differ, the mixed-state comparison simply chooses the better orientation.
This is useful: the same neighborhood test can support compression
whether the current gaps are small, large, or asymmetric.

\section{The resulting class and the remaining dichotomy}

The class can be dense and non-bipartite. Start with arbitrary quadratic
interactions on $n/2$ classes, replace each class by a pair, and make
interactions constant across each pair of classes after optional variable
negations. Internal pair interactions and all fields are arbitrary.
Perturb a bounded number of coefficients incident to each pair. Taking
both internal and class-to-class interactions nonzero gives a complete interaction
graph, with density proportional to $n$. The quotient can contain a
linear number of classes; there need not be a small separator or a small
total number of neighborhood types. This is more permissive than
requiring a few exact interaction classes.

The theorem establishes one structural exit when neighborhood responses
are nearly identical. It does \emph{not} establish the converse:
frequent small total gap need not yet be known to force linearly many
such pairs. The remaining inverse question is whether persistent small
gap yields near-twin pairs, larger compressible groups, a small set of
shared influential variables, or some other structure with a provable
algorithmic saving.

For integration into a pruning search, the theorem is directly safe as
a terminal solver of the residual \emph{soft} objective. Auxiliary
pruning constraints may break the pair symmetry; discarding them for
this terminal optimization is safe because every returned assignment is
still valid for the original soft problem. Genuinely imposed hard
constraints cannot be discarded this way.

\begingroup\small\raggedright
\paragraph{Relation to the previous notes.}
This develops the approximate interaction-class direction left open in
\texttt{max2sat\_branching\_frontier.tex} and the pair-decision direction
in \texttt{max2sat\_pairs\_sums\_measure.tex}. The agreement counts of
\texttt{max2sat\_counting\_conditions.tex} concern cancellation; the
near-twin test here instead detects strong alignment that can be
exploited by domination after resolving discrepancies.

\paragraph{Verification.}
The companion program \texttt{explore\_near\_twin\_pairs.py} compares
the three-state algorithm with exhaustive optimization on dense planted
examples with negated pairs, unequal fields, exceptional edges, and
unpaired variables. All 26 instances preserved the optimum, comparing
21,795 compressed candidates with 59,648 exhaustive assignments.
There were 2,160 gap-identity checks, and 30 retained pairs whose chosen
mixed state changed between outside assignments. The compatibility
lemmas passed checks on all 324 oriented $2\times2$ interaction blocks
over $\{-1,0,+1\}$ and all 1,099 labeled graphs on one to five vertices.
Results are recorded in \texttt{near\_twin\_pair\_checks.json}.
\par\endgroup
\end{document}
